% Additive manuscript section. Dependencies: amsmath, amssymb, amsthm;
% proposition and proof environments. Bibliography items supplied separately.
\section{Centered Rogers--Dougall resonance}
\label{dr-core:section}

Fix $\kappa,b,c>0$ with $d_0=1+2\kappa-b-c>0$, and write $x_n=n+\kappa$.
For complex $u$ with $\operatorname{Re}u<d_0$, put
\begin{align}
 W_n(u)={}&x_n\frac{\Gamma(n+2\kappa)\Gamma(n+b)}
 {\Gamma(n+1)\Gamma(n+1+2\kappa-b)}
 \frac{\Gamma(n+c)\Gamma(n+d_0-u)}
 {\Gamma(n+1+2\kappa-c)\Gamma(n+b+c+u)},
 \label{dr-core:weight}\\
 P(u)={}&\frac{\Gamma(b)\Gamma(c)\Gamma(d_0-u)}
 {\Gamma(d_0)\Gamma(b+u)\Gamma(c+u)}.
 \label{dr-core:P}
\end{align}
Denominator Gamma factors always mean entire reciprocal-Gamma factors.
The classical Rogers--Dougall sum gives
$\sum_{n\geq0}W_n(u)=\Gamma(u)P(u)/2$ for $0<\operatorname{Re}u<d_0$
\cite{dougall-resonance:DLMF16}. Its hypergeometric excess is $2u$.

Set $r=(\kappa,b-\kappa,c-\kappa,1+\kappa-b-c-u)$ and define
\begin{align}
 \ell_j(u)&=-\frac{\sum_{i=1}^4B_{2j+1}(r_i)}{j(2j+1)},\qquad j\geq1,
 \notag\\
 A_0(u)&=1,\qquad
 A_j(u)=\frac1j\sum_{h=1}^j h\ell_h(u)A_{j-h}(u).
 \label{dr-core:A}
\end{align}
The centered Gamma form is
$W_n(u)=x_n\prod_{i=1}^4\Gamma(x_n+r_i)/\Gamma(x_n+1-r_i)$.
Bernoulli reflection and the shifted Stirling expansion imply
\begin{equation}
 W_n(u)=x_n^{-1-2u}\left(\sum_{j=0}^{K-1}A_j(u)x_n^{-2j}
                                    +O(x_n^{-2K})\right).
 \label{dr-core:asymptotic}
\end{equation}
Only even inverse-power corrections occur. The remainder is locally
uniform in the parameters; fixed spectral derivatives introduce only
powers of $\log x_n$ \cite{dougall-resonance:DLMF5}.

\begin{proposition}[An ordinary convergent subtraction]
\label{dr-core:subtraction}
For $K\geq1$ and $-K<\operatorname{Re}u<d_0$,
\begin{align}
 \mathcal R_K(u)
 &:={\sum_{n=0}^{\infty}}
 \left[W_n(u)-\sum_{j=0}^{K-1}A_j(u)x_n^{-1-2u-2j}\right]\notag\\
 &=\frac12\Gamma(u)P(u)-\sum_{j=0}^{K-1}
                              A_j(u)\zeta(1+2u+2j,\kappa).
 \label{dr-core:R}
\end{align}
The bracketed series and every fixed spectral derivative converge normally
in that strip. Apparent poles on the right side are removable.
\end{proposition}
\begin{proof}
Equation~\eqref{dr-core:asymptotic} gives a summable compact-uniform
majorant $O(n^{-1-\epsilon}(1+\log^m n))$ for each fixed derivative.
On $0<\operatorname{Re}u<d_0$, sum the components separately by classical
Dougall summation. The identity theorem then gives \eqref{dr-core:R}
on the full strip, with removability supplied by its holomorphic left side.
\end{proof}

\begin{proposition}[Residue polynomial and complete collapse]
\label{dr-core:collapse}
Let $N\geq0$ and
\begin{equation}
 \rho_N=\frac{(-1)^N}{N!}(d_0)_N(1-b)_N(1-c)_N.
 \label{dr-core:rho}
\end{equation}
Then $A_N(-N)=\rho_N$, and
\begin{equation}
 \mathcal R_{N+1}(-N)=\frac{\rho_N}{2}
 [\psi(N+1)-\psi(d_0+N)-\psi(b)-\psi(c)+2\psi(\kappa)].
 \label{dr-core:value}
\end{equation}
All displayed digamma arguments are positive. In particular, the formula
also applies when $\rho_N=0$.
\end{proposition}
\begin{proof}
Compare residues in \eqref{dr-core:R}. The Gamma term has residue
$\rho_N/2$ at $u=-N$, and the only polar counterterm is $j=N$, with
residue $A_N(-N)/2$. This proves the polynomial identity $A_N(-N)=\rho_N$.

The essential coefficient law is
\begin{equation}
 (\partial_u-D)A_N(u)=\sum_{h=1}^N\frac{B_{2h}(\kappa)}hA_{N-h}(u),
 \qquad D=\partial_\kappa+\partial_b+\partial_c.
 \label{dr-core:law}
\end{equation}
Indeed, $Dd_0=0$ and $E=\partial_u-D$ acts on the four $r_i$ by
$Er_1=-1$ and $Er_2=Er_3=Er_4=0$. Hence
$E\ell_h=B_{2h}(\kappa)/h$. Apply $E$ to the formal generating function
$\sum A_jy^j=\exp(\sum\ell_jy^j)$ and extract $y^N$.

For generic $b,c$, the constant term of $\Gamma(-N+t)P(-N+t)/2$ is
\[
 \frac{\rho_N}{2}
 [\psi(N+1)-\psi(d_0+N)-\psi(b-N)-\psi(c-N)].
\]
The $j=N$ counterterm has constant term
$A_N'(-N)/2-\rho_N\psi(\kappa)$, while the lower counterterms sum to
$-\sum_{h=1}^N B_{2h}(\kappa)A_{N-h}(-N)/(2h)$.
Equation~\eqref{dr-core:law} cancels all of the latter terms, leaving
$-D\rho_N/2$ in place of the coefficient derivative. Finally,
\[
 D\rho_N=\rho_N\sum_{h=1}^N\bigl((b-h)^{-1}+(c-h)^{-1}\bigr).
\]
The digamma shift relation proves \eqref{dr-core:value}. Continuity of
the normally convergent series and of the final positive-argument formula
proves the exceptional cases as well.
\end{proof}

The full spectral germ retains reciprocal-Gamma zeros. Define
\begin{equation}
 Q_N(t)=\frac{(-1)^N}{N!}\Gamma(1+t)
             \prod_{h=1}^N(1-t/h)^{-1}P(-N+t).
 \label{dr-core:Q}
\end{equation}
It is analytic at zero, $Q_N(0)=\rho_N$, and
$\Gamma(-N+t)P(-N+t)/2=Q_N(t)/(2t)$.
With $a_j(t)=A_j(-N+t)$ and
$\zeta(1+v,\kappa)=v^{-1}+\sum_{m\geq0}(-1)^m\gamma_m(\kappa)v^m/m!$,
\begin{equation}
\begin{split}
 \mathcal R_{N+1}(-N+t)={}&\frac{Q_N(t)-a_N(t)}{2t}
 -a_N(t)\sum_{m\geq0}\frac{(-2)^m\gamma_m(\kappa)}{m!}t^m\\
 &-\sum_{j=0}^{N-1}a_j(t)\zeta(1-2(N-j)+2t,\kappa).
\end{split}
\label{dr-core:jet}
\end{equation}
This is a holomorphic identity and therefore proves every Taylor
coefficient formula. Only negative odd Hurwitz spectral arguments occur
outside the Stieltjes part. To evaluate at a zero of a reciprocal Gamma
factor, use
\[
 \frac1{\Gamma(w+t)}=
 \frac{\prod_{h=0}^{L-1}(w+h+t)}{\Gamma(w+L+t)},\qquad w+L>0,
\]
before expanding. A vanished summand is not necessarily a vanished jet.

At the interlaced point $u_*=-N-1/2$, ordinary substitution instead gives
\begin{equation}
 \mathcal R_{N+1}(u_*)=\frac12\Gamma(u_*)P(u_*)+
 \sum_{j=0}^N A_j(u_*)\frac{B_{2(N-j)+1}(\kappa)}{2(N-j)+1}.
 \label{dr-core:half}
\end{equation}
Its spectral jets involve $0,-2,\ldots,-2N$, with no Stieltjes pole.
These conclusions follow from the classical Bernoulli special values of
Hurwitz zeta \cite{dougall-resonance:DLMF25} and \eqref{dr-core:R}.

For $|z|<1$, replace each counterterm by
$A_j(u)\Phi(z,1+2u+2j,\kappa)$ and $W_n(u)$ by $W_n(u)z^n$.
The same bracketed convergence theorem extends all fixed spectral jets
continuously to $z=1$. If $\kappa=p/q$, $1\leq p\leq q$, then
\[
 \Phi(z,s,p/q)=q^{s-1}z^{-p/q}\sum_{\ell=0}^{q-1}
       e^{-2\pi i p\ell/q}\operatorname{Li}_s(e^{2\pi i\ell/q}z^{1/q}),
\]
initially for $0<z<1$. This gives exact polylogarithmic subtraction
identities. At the endpoint, take the limit of the complete bracket,
not of separately divergent pieces.

The classical Dougall sum is the input, not a new claim. The centered
coefficient law and the explicit resonance/jet calculus are the continuation
presented here. The larger article supplies detailed harmonic and primitive
specializations, antecedent discussion, and reproducible diagnostics.
